\chapter{Wyniki i analiza}
\label{cha:elementyPracy}

\section{Dokładność na zbiorze testowym}

\begin{table}[H]
    \centering
    \caption{Porównanie dokładności klasyfikatorów na zbiorze testowym Sign Language MNIST}
    \label{tab:wyniki}
    \begin{tabular}{llr}
        \toprule
        \textbf{Model} & \textbf{Plik artefaktu} & \textbf{Dok. testowa} \\
        \midrule
        SVC (jądro RBF)        & \texttt{artifacts/svc\_model.pkl}   & 0{,}82 \\
        Las Losowy (200 drzew) & \texttt{artifacts/rf\_model.pkl}    & 0{,}82 \\
        CNN (Keras/TF)         & \texttt{artifacts/cnn\_model.keras} & \textbf{0{,}92} \\
        \bottomrule
    \end{tabular}
\end{table}

CNN osiąga wyraźnie wyższą dokładność (92\%) w porównaniu z SVC i RF (oba~82\%).
Wynik jest spójny z literaturą --- architektury konwolucyjne są naturalnie
przystosowane do danych obrazowych.

\section{Dyskusja wyników}

\subsection{Przewaga CNN nad metodami klasycznymi}

Różnica 10 punktów procentowych na korzyść CNN wynika z fundamentalnej cechy
architektury: warstwy konwolucyjne uczą się hierarchicznych reprezentacji
przestrzennych --- krawędzi, narożników i wzorców wyższego rzędu --- których
płaski wektor 784 cech nie przekazuje. SVC i RF operują na przestrzeni
784-wymiarowej bez żadnego uwzględnienia sąsiedztwa pikseli.

\subsection{Równoważność SVC i Lasu Losowego}

Zbliżona dokładność obu metod (82\%) sugeruje, że granica decyzyjna
jest podobnie złożona: wystarczająco skomplikowana, by wymagać nieliniowych
metod, ale niewystarczająco złożona, by jedna metoda miała wyraźną przewagę
bez dalszego strojenia hiperparametrów.

\section{Ograniczenia obecnej wersji}
\label{sec:ograniczenia}

Wersja aktualna jest rozwiązaniem klasyfikacji pojedynczych klatek i nie
modeluje dynamiki ruchu. Oznacza to, że:
\begin{itemize}
    \item \textbf{Znak J i Z} --- znaki wymagające ruchu (np. J, Z) nie są klasyfikowane,
    \item \textbf{Brak mechanizmu} ---  wygładzania predykcji po czasie,
    \item \textbf{Brak detekcji kontekstu} --- kontekstu sekwencji gestów.
    \item \textbf{Wrażliwość na oświetlenie} --- zmiana natężenia lub kierunku
    światła modyfikuje wartości pikseli i może prowadzić do błędnych
    predykcji.
    \item \textbf{Wrażliwość na pozycję} --- dłoń musi być wycentrowana w kadrze
        i podobnej skali co próbki treningowe.
    \item \textbf{Brak niezmienności geometrycznej} --- modele nie są z natury
    niezmiennicze na obrót, skalę ani przesunięcie; CNN częściowo
    kompensuje to przez operację \textit{max pooling}.
\end{itemize}

Mimo tych ograniczeń projekt spełnia cel dydaktyczny i demonstracyjny:
spójny pipeline od danych tablicowych do inferencji online.
